\chapter{Seguridad}
\label{cap:seguridad}

Este cap\'itulo integra la auditor\'ia de seguridad de BIOPET desde tres
\'angulos complementarios: revisi\'on manual sobre OWASP Top~10~\cite{owasp-top10},
an\'alisis din\'amico (OWASP ZAP Baseline) y an\'alisis est\'atico
(SpotBugs + Find Security Bugs), m\'as la auditor\'ia dedicada de SQL
din\'amico en procedimientos almacenados.

\section{Autenticaci\'on por cookies y revocaci\'on}

JWT firmado con HMAC-SHA256, transportado en cookies
\texttt{HttpOnly}+\texttt{Secure}+\texttt{SameSite=Strict}
(\texttt{access\_token}, \texttt{refresh\_token}), con revocaci\'on real
v\'ia lista negra en Redis indexada por \texttt{jti} (ADR-003), rate
limiting de login en memoria (5 fallos consecutivos $\to$ 401; el sexto
$\to$ 429 con cabecera \texttt{Retry-After}) y auditor\'ia estructurada
(\texttt{AUTH\_AUDIT}) que solo acepta \texttt{ip} y \texttt{subject} como
par\'ametros, sin v\'ia de c\'odigo para que una contrase\~na, un JWT
completo o un encabezado \texttt{Authorization} llegue al logger.

\section{Autorizaci\'on por rol y por propiedad}

Autorizaci\'on combinada: por rol (\texttt{@PreAuthorize}) y por propiedad
(\texttt{ROLE\_DUENO} solo accede a sus propias mascotas). 401 para
autenticaci\'on ausente/inv\'alida, 403 para autenticaci\'on v\'alida sin
autorizaci\'on, ambos como \texttt{ProblemDetail} sin \emph{stack trace}
(figura~\ref{fig:postman-401-final}).

% EVIDENCIA-JAIME-05
% Ruta esperada: figuras/jaime/05-postman-401.png
\evidencia{figuras/jaime/05-postman-401.png}%
{Respuesta 401 (\texttt{ProblemDetail}) capturada desde la colecci\'on Postman de BIOPET.}%
{fig:postman-401-final}

\section{Cabeceras de seguridad y TLS}

\texttt{SecurityConfig} (Spring Security~\cite{spring-security-docs}) fija
cinco cabeceras en todas las respuestas: \texttt{X-Frame-Options: DENY},
\texttt{X-Content-Type-Options: nosniff},
\texttt{Content-Security-Policy: default-src 'self'; frame-ancestors 'none'; object-src 'none'},
\texttt{Referrer-Policy: no-referrer}, y \texttt{Strict-Transport-Security}
solo sobre HTTPS. TLS~1.3 verificado en vivo
(\texttt{TLS\_AES\_256\_GCM\_SHA384}, figura~\ref{fig:tls-openssl-final})
en el entorno de desarrollo; en
producci\'on, HTTPS es provisto autom\'aticamente por Render sobre el
subdominio \texttt{*.onrender.com} (cap\'itulo~\ref{cap:despliegue}).

% EVIDENCIA-JAIME-03
% Ruta esperada: figuras/jaime/03-tls-openssl.png
\evidencia{figuras/jaime/03-tls-openssl.png}%
{Salida de \texttt{openssl s\_client -tls1\_3}: TLSv1.3 y \texttt{TLS\_AES\_256\_GCM\_SHA384} negociados (entorno de desarrollo).}%
{fig:tls-openssl-final}

% EVIDENCIA-JAIME-04
% Ruta esperada: figuras/jaime/04-security-headers.png
\evidencia{figuras/jaime/04-security-headers.png}%
{Cabeceras de seguridad reales capturadas con \texttt{curl.exe -i}.}%
{fig:security-headers-final}

\section{Auditor\'ia OWASP Top~10 (resumen)}

La tabla~\ref{tab:owasp-final} resume el resultado por control de la
revisi\'on manual sobre los nueve controles OWASP Top~10 con evidencia
real y verificable.

\begin{table}[H]
  \centering
  \caption{Resultado de la auditor\'ia manual OWASP por control (fuente: \rutalarga{docs/mediciones/sec/REPORT.md}).}
  \label{tab:owasp-final}
  \begin{tabular}{@{}ll@{}}
    \toprule
    Control OWASP & Resultado \\
    \midrule
    A01 --- Control de acceso roto & PASS \\
    A02 --- Fallas criptogr\'aficas (TLS) & PASS \\
    A03 --- Inyecci\'on & PASS \\
    A04 --- Dise\~no inseguro & Evaluado (cierre 2026-08-09) \\
    A05 --- Configuraci\'on de seguridad & PASS \\
    A06 --- Componentes vulnerables/desactualizados & Evaluado, con hallazgos documentados \\
    A07 --- Identificaci\'on y autenticaci\'on & PASS \\
    A09 --- Registro y monitoreo & PASS \\
    XSS (transversal) & Evaluado (cierre 2026-08-09) \\
    \bottomrule
  \end{tabular}
\end{table}

A08 y A10 no fueron objeto de auditor\'ia dedicada y no se documentan aqu\'i
para no inventar evidencia inexistente.

\section{Protecci\'on estructural contra inyecci\'on (A03)}

Todos los m\'etodos de repositorio son consultas derivadas de Spring Data o
invocaciones formales de procedimientos almacenados
(\texttt{@Procedure}/\texttt{@NamedStoredProcedureQuery}), con par\'ametros
enlazados y fuertemente tipados. Los payloads de inyecci\'on probados
(\texttt{1 OR 1=1}, \texttt{1; DROP TABLE usuarios; --}, \texttt{\%' UNION
SELECT NULL --}) son rechazados en la capa de conversi\'on de Spring MVC
antes de alcanzar el repositorio, verificado en \texttt{SqlInjectionSecurityTest}.

\section{Auditor\'ia est\'atica: SpotBugs + Find Security Bugs}
\label{sec:spotbugs}

Ejecutada con \texttt{spotbugs-maven-plugin}~4.10.3.0 + Find Security
Bugs~1.14.0, \texttt{effort=Max}, \texttt{threshold=Low} (m\'axima
exhaustividad), sobre todas las clases compiladas del backend; la
tabla~\ref{tab:spotbugs-final} resume los hallazgos por severidad.

\begin{table}[H]
  \centering
  \caption{Resumen del an\'alisis est\'atico (fuente: \rutalarga{docs/mediciones/sec/static-analysis/README.md}).}
  \label{tab:spotbugs-final}
  \begin{tabular}{@{}lr@{}}
    \toprule
    M\'etrica & Valor \\
    \midrule
    Hallazgos totales & 66 \\
    Severidad alta (\texttt{priority=1}) & 1 (\texttt{SPRING\_CSRF\_PROTECTION\_DISABLED}) \\
    Severidad media (\texttt{priority=2}) & 13 \\
    Severidad baja (\texttt{priority=3}) & 52 \\
    Hallazgos relacionados con SQL (\texttt{SQL\_*}) & \textbf{0} \\
    \bottomrule
  \end{tabular}
\end{table}

\subsection*{Tratamiento del hallazgo CSRF (\'unico de severidad alta)}
\label{sec:csrf}

\texttt{SecurityConfig.java} deshabilita expl\'icitamente la protecci\'on
CSRF integrada de Spring Security, lo que SpotBugs/Find Security Bugs
reporta correctamente como \texttt{SPRING\_CSRF\_PROTECTION\_DISABLED}
(\texttt{priority=1}). Este informe \textbf{no trata este hallazgo como un
falso positivo}: la herramienta detect\'o exactamente lo que hay en el
c\'odigo. Se trata, en cambio, como un \textbf{riesgo mitigado por
arquitectura}: la protecci\'on CSRF tradicional basada en tokens existe
para defender flujos de autenticaci\'on por \emph{sesi\'on con cookie
ambiental} en formularios HTML cl\'asicos, donde el navegador env\'ia la
cookie de sesi\'on autom\'aticamente en cualquier solicitud
\emph{cross-site}, incluida una iniciada por un sitio malicioso. BIOPET no
sigue ese patr\'on: el atributo \texttt{SameSite=Strict} en las cookies
\texttt{access\_token}/\texttt{refresh\_token} (ADR-006) impide que el
navegador adjunte esas cookies en solicitudes iniciadas desde un origen
distinto, que es precisamente el vector que CSRF explota. La combinaci\'on
\texttt{HttpOnly}+\texttt{Secure}+\texttt{SameSite=Strict} sustituye
funcionalmente al token CSRF cl\'asico para este esquema de autenticaci\'on
espec\'ifico, sin que eso implique que el riesgo no exist\'ia --- exist\'ia,
y se document\'o expl\'icitamente en vez de suprimirse. El gate de CI
\texttt{security-static} (secci\'on~\ref{sec:ci-jobs}) falla \'unicamente
ante hallazgos \texttt{SQL\_*}; el hallazgo CSRF se archiva siempre como
parte del reporte completo (sin filtros ni \texttt{@SuppressFBWarnings}),
para que quede visible en cada corrida en vez de ocultarse.

\section{Auditor\'ia din\'amica: OWASP ZAP}
\label{sec:zap}

Esta secci\'on documenta \textbf{dos} corridas de OWASP ZAP, claramente
diferenciadas: (A) el \emph{baseline} hist\'orico sin autenticaci\'on
(evidencia original de esta fase, conservada sin alterar), y (B) un
escaneo pasivo \textbf{autenticado}, incorporado en la recalificaci\'on
para cerrar la limitaci\'on que (A) ya declaraba expl\'icitamente. Ninguna
de las dos corridas es un \emph{active scan}: ambas son \emph{baseline}/pasivas,
sin intentos de explotaci\'on.

\subsection{A. Baseline hist\'orico (sin autenticaci\'on)}
\label{sec:zap-baseline-historico}

Ejecutada con \texttt{zap-baseline.py} (imagen oficial
\texttt{ghcr.io/zaproxy/zaproxy:stable}, versi\'on~2.17.0) contra el
backend real, dentro de la red Docker Compose (\texttt{http://backend:8080},
\textbf{sin TLS}, \textbf{sin autenticaci\'on}), el
\texttt{2026-08-18T02:53:21Z} (fuente:
\texttt{docs/mediciones/sec/zap/RUN-METADATA.txt}).

\begin{table}[H]
  \centering
  \caption{Resumen ZAP Baseline hist\'orico (fuente: \texttt{docs/mediciones/sec/zap/README.md}).}
  \label{tab:zap-final}
  \begin{tabular}{@{}lr@{}}
    \toprule
    Severidad & Alertas \\
    \midrule
    Alta & \textbf{0} \\
    Media & 0 \\
    Baja & 0 \\
    Informativa & 1 (\texttt{Non-Storable Content}, 3 instancias) \\
    \bottomrule
  \end{tabular}
\end{table}

C\'odigo de salida de \texttt{zap-baseline.py}: \texttt{0}. La
tabla~\ref{tab:zap-final} resume la severidad de las alertas; en t\'erminos
de cobertura real, el \emph{spider} no autenticado solo alcanz\'o
\textbf{3~URI} (\texttt{/}, \texttt{/robots.txt}, \texttt{/sitemap.xml} ---
las tres instancias de \texttt{Non-Storable Content} de la tabla) sobre
\textbf{2~endpoints \'unicos} seg\'un \texttt{insight.endpoint.total} del
propio JSON de ZAP (m\'etrica de ZAP que deduplica por endpoint, distinta
del conteo de URI), con \textbf{100\,\%} de las respuestas en \texttt{4xx}
(\texttt{insight.code.4xx}). El objetivo de la Entrega Final (cero
hallazgos de severidad alta) se cumple, pero esta corrida \textbf{no
constituye evidencia representativa de la superficie protegida}: el
spider de ZAP no fue autenticado y, por dise\~no, todo lo que no sea
\texttt{/api/auth/**} responde \texttt{401} (limitaci\'on declarada
expl\'icitamente en \texttt{docs/mediciones/sec/zap/README.md}, no
oculta).

\subsection{B. Recalificaci\'on --- escaneo pasivo autenticado (local, TLS)}
\label{sec:zap-authenticated-local}

Para cerrar la limitaci\'on de (A), se ejecut\'o un segundo escaneo con
OWASP ZAP~2.17.0 (misma imagen), esta vez \textbf{autenticado} y sobre
\textbf{HTTPS/TLS} real, contra el mismo backend pero levantado con el
perfil TLS local del propio repositorio
(\texttt{docker-compose.tls.yml}, \texttt{https://backend:8443}), el
\texttt{2026-09-03T21:38:48Z} (fuente:
\texttt{docs/mediciones/sec/zap/RUN-METADATA-AUTH-LOCAL.txt}). Sigue
siendo un entorno \textbf{local}, del propio repositorio acad\'emico ---
\textbf{no} se escane\'o el despliegue de Render ni ning\'un otro
repositorio (BIOPET-V2 es un proyecto separado, fuera de alcance).

La autenticaci\'on us\'o el m\'etodo \texttt{json} del \emph{Automation
Framework} de ZAP contra \texttt{POST~/api/auth/login}, con gesti\'on de
sesi\'on por \textbf{cookie} (m\'etodo nativo de ZAP, coincide con el
mecanismo real de \texttt{JwtCookieService}) y verificaci\'on peri\'odica
contra \texttt{GET~/api/usuarios/me}. El login y las siete rutas GET
protegidas objetivo (\texttt{/api/usuarios/me}, \texttt{/api/mascotas},
\texttt{/api/mascotas/resumen-especies}, \texttt{/api/citas},
\texttt{/api/consultas}, \texttt{/api/vacunas}, \texttt{/api/usuarios})
devolvieron exactamente \texttt{200}, verificado con la aserci\'on nativa
\texttt{responseCode: 200} de ZAP AF en cada petici\'on, con \textbf{0
diferencias} registradas.

Seg\'un los \emph{insights} agregados del propio JSON de ZAP,
\texttt{insight.endpoint.total = 9} (endpoints \'unicos, deduplicados);
el resumen de consola de \texttt{outputSummary} reporta, por separado,
``Total of 11 URLs'' --- un contador \textbf{distinto}, sin deduplicar
(incluye la repetici\'on de \texttt{POST~/api/auth/login}, invocado dos
veces por dise\~no del ciclo de verificaci\'on). \textbf{Estas dos cifras
no deben sumarse ni tratarse como la misma m\'etrica.} De igual modo,
\texttt{insight.code.2xx = 93} e \texttt{insight.code.4xx = 6} son
estad\'isticas agregadas internas de ZAP, \textbf{no} un desglose por
URL; ninguno de los reportes retenidos (HTML, XML ni JSON) expone el
denominador exacto de ese porcentaje, por lo que este informe
\textbf{no afirma} ``93\,\% de 11 peticiones'' ni ninguna fracci\'on
literal equivalente. Lo que s\'i qued\'o verificado por partida doble es
que ese \texttt{4xx} agregado \textbf{no corresponde a ninguna de las
siete rutas objetivo ni al login} (ambos confirmados en \texttt{200} por
la aserci\'on \texttt{responseCode} con 0 diferencias) y que una
reconciliaci\'on adicional de solo lectura, exportando el historial HTTP
real de una corrida equivalente, mostr\'o \textbf{0 respuestas 4xx} entre
las nueve peticiones expl\'icitas del plan. La identidad exacta de la
petici\'on individual que origin\'o ese \texttt{4xx} agregado no pudo
reconstruirse a partir de los artefactos retenidos --- se documenta aqu\'i
como \textbf{l\'imite de trazabilidad} de las plantillas de reporte de
ZAP, no como un hecho supuesto.

Alertas de esta corrida: \textbf{0} altas, \textbf{0} medias, \textbf{0}
bajas, \textbf{2 tipos informativos} (\texttt{Authentication Request
Identified} y \texttt{Session Management Response Identified}, 3
instancias en total, ambas sobre \texttt{POST~/api/auth/login}) --- son
ZAP reconociendo correctamente el propio flujo de autenticaci\'on, \textbf{no}
vulnerabilidades.

La tabla~\ref{tab:zap-comparativo} resume la comparaci\'on entre ambas
corridas.

\begin{table}[H]
  \centering
  \caption{Comparaci\'on entre el baseline hist\'orico sin autenticar y el escaneo pasivo autenticado local (TLS).}
  \label{tab:zap-comparativo}
  \begin{tabular}{@{}lll@{}}
    \toprule
    Criterio & A. Baseline hist\'orico & B. Autenticado local (TLS) \\
    \midrule
    Target & \texttt{http://backend:8080} & \texttt{https://backend:8443} \\
    TLS & No & S\'i (autofirmado, local) \\
    Autenticaci\'on & No & S\'i (\texttt{json}, cookie) \\
    Rutas protegidas en \texttt{2xx} & 0 & 7/7 \\
    Endpoints \'unicos & 2 & 9 \\
    Alta & 0 & 0 \\
    Media & 0 & 0 \\
    Baja & 0 & 0 \\
    Limitaci\'on principal & Sin cobertura autenticada & Solo GET, solo local \\
    \bottomrule
  \end{tabular}
\end{table}

\textbf{Alcance expl\'icito de esta corrida (B) --- lo que NO se hizo:}
no se ejecut\'o ning\'un \emph{active scan}; no se atacaron rutas
\texttt{POST}/\texttt{PUT}/\texttt{DELETE} (solo las siete rutas
\texttt{GET} listadas); no se cubri\'o el 100\,\% de los endpoints de la
API (quedan fuera, deliberadamente, todas las rutas de escritura y
\texttt{/api/externa/**}, un proxy a una API externa real); no se
escane\'o el despliegue de Render ni producci\'on; no se escane\'o
BIOPET-V2. Esta corrida \textbf{no valida la seguridad completa del
sistema}; mejora sustancialmente la representatividad del DAST porque
verifica autenticaci\'on real y rutas protegidas reales, sobre un
subconjunto acotado y deliberado de la superficie de la API. Extender
este mismo mecanismo de autenticaci\'on a un escaneo m\'as amplio, o
contra un entorno p\'ublico/de producci\'on, es una extensi\'on posible
para trabajo futuro que requerir\'ia adem\'as control expl\'icito de
credenciales y de la carga generada contra un servicio compartido.

Evidencia completa, trazable y verificable por SHA-256 en
\texttt{docs/mediciones/sec/zap/} (\texttt{zap-authenticated-local-report.json},
\texttt{.html}, \texttt{.xml}, \texttt{RUN-METADATA-AUTH-LOCAL.txt},
\texttt{zap-authenticated-local.yaml}, \texttt{SHA256SUMS-AUTH-LOCAL.txt})
y en \texttt{scripts/run-zap-authenticated-local.sh} (script reproducible,
sin credenciales hardcodeadas). Verificaci\'on:
\texttt{sha256sum -c SHA256SUMS-AUTH-LOCAL.txt} $\to$ \textbf{7/7 OK}.

\section{Auditor\'ia de SQL din\'amico en procedimientos almacenados}

\texttt{scripts/audit-sql-dynamic.sh} analiza los seis archivos de
\texttt{db/procs/*.sql} en busca de construcci\'on insegura de SQL
din\'amico (\texttt{EXECUTE} con concatenaci\'on, patrones ajenos a
PostgreSQL). Resultado: \textbf{0 hallazgos}, \texttt{exit 0}.

\section{Postman}

La colecci\'on \rutalarga{docs/postman/BIOPET.postman_collection.json}
(40+ \emph{requests}) y \rutalarga{docs/postman/BIOPET-Vacunas.postman_collection.json}
cubren estado del servicio, autenticaci\'on completa (incluida la
secuencia de rate limiting), mascotas, citas, consultas, vacunas y un
cat\'alogo de errores controlados (422, 400, 404, 401, 403, 429). Ninguna
depende de un JWT o cookie guardado manualmente en variable: se apoyan
exclusivamente en el \emph{cookie jar} autom\'atico de Postman.

\section{Limitaciones reales de este cap\'itulo}

\begin{itemize}[nosep]
  \item El certificado TLS del entorno de desarrollo es autofirmado y
    exclusivamente acad\'emico; en producci\'on, HTTPS lo provee Render de
    forma autom\'atica (no es un certificado gestionado por el equipo).
  \item El rate limiting de login es en memoria y por instancia: no es
    distribuido entre m\'ultiples r\'eplicas del backend.
  \item Los logs \texttt{AUTH\_AUDIT} se escriben \'unicamente en el log
    local del proceso; no existe integraci\'on con un SIEM centralizado.
  \item El \emph{spider} del baseline hist\'orico de ZAP
    (secci\'on~\ref{sec:zap-baseline-historico}) no fue autenticado; esa
    limitaci\'on espec\'ifica qued\'o parcialmente cerrada por el escaneo
    pasivo autenticado local (secci\'on~\ref{sec:zap-authenticated-local}),
    que verific\'o en \texttt{200} siete rutas \texttt{GET} protegidas
    reales de los m\'odulos de mascotas, citas, consultas, vacunas y
    usuarios --- pero sigue sin cubrir las rutas de escritura
    (\texttt{POST}/\texttt{PUT}/\texttt{DELETE}) de esos mismos m\'odulos,
    ni el despliegue de Render, ni un \emph{active scan}.
  \item El an\'alisis est\'atico no sustituye una revisi\'on de seguridad
    din\'amica dirigida a la librer\'ia JWT subyacente (\texttt{jjwt}); ese
    tipo de an\'alisis (\emph{fuzzing} dirigido, ver
    Yang et al.~\cite{yang2026tokentimebomb}, cap\'itulo~\ref{cap:trabajos-relacionados})
    queda fuera del alcance de esta auditor\'ia.
  \item ``0 alertas de severidad alta''/``0 hallazgos \texttt{SQL\_*}'' no
    implica ausencia total de vulnerabilidades: implica ausencia de los
    patrones que las reglas configuradas de ZAP y SpotBugs pueden detectar.
\end{itemize}
